Les pistes d'un disque optique sont séparées de \qty{\v{pas}}{\micro\metre} et se comportent comme un réseau par réflexion. On envoie perpendiculairement sur le disque le faisceau d'un pointeur laser de longueur d'onde \qty{\v{lam}}{nm}, et on observe des taches lumineuses sur un mur.

\begin{enumerate}[label=\alph*)]
    \item À quel angle de la direction de réflexion directe se trouve le maximum d'ordre 1 ?
    \item Quel est l'ordre le plus élevé observable ?
    \item Combien de taches lumineuses au total (en comptant la tache centrale) peut-on observer ?
\end{enumerate}
